% Options for packages loaded elsewhere
\PassOptionsToPackage{unicode}{hyperref}
\PassOptionsToPackage{hyphens}{url}
%
\documentclass[
  11pt,
]{ctexart}
\usepackage{amsmath,amssymb}
\usepackage{iftex}
\ifPDFTeX
  \usepackage[T1]{fontenc}
  \usepackage[utf8]{inputenc}
  \usepackage{textcomp} % provide euro and other symbols
\else % if luatex or xetex
  \usepackage{unicode-math} % this also loads fontspec
  \defaultfontfeatures{Scale=MatchLowercase}
  \defaultfontfeatures[\rmfamily]{Ligatures=TeX,Scale=1}
\fi
\usepackage{lmodern}
\ifPDFTeX\else
  % xetex/luatex font selection
\fi
% Use upquote if available, for straight quotes in verbatim environments
\IfFileExists{upquote.sty}{\usepackage{upquote}}{}
\IfFileExists{microtype.sty}{% use microtype if available
  \usepackage[]{microtype}
  \UseMicrotypeSet[protrusion]{basicmath} % disable protrusion for tt fonts
}{}
\makeatletter
\@ifundefined{KOMAClassName}{% if non-KOMA class
  \IfFileExists{parskip.sty}{%
    \usepackage{parskip}
  }{% else
    \setlength{\parindent}{0pt}
    \setlength{\parskip}{6pt plus 2pt minus 1pt}}
}{% if KOMA class
  \KOMAoptions{parskip=half}}
\makeatother
\usepackage{xcolor}
\usepackage[top=23mm,bottom=23mm,left=24mm,right=24mm]{geometry}
\setlength{\emergencystretch}{3em} % prevent overfull lines
\providecommand{\tightlist}{%
  \setlength{\itemsep}{0pt}\setlength{\parskip}{0pt}}
\setcounter{secnumdepth}{-\maxdimen} % remove section numbering
\usepackage{fontspec}
\usepackage{xeCJK}
\setmainfont{Noto Serif}
\setsansfont{Noto Sans}
\setmonofont{Noto Sans Mono}
\setCJKmainfont[ItalicFont={Noto Serif CJK SC}]{Noto Serif CJK SC}
\setCJKsansfont{Noto Sans CJK SC}
\usepackage{microtype}
\usepackage{setspace}
\setstretch{1.20}
\usepackage[most]{tcolorbox}
\usepackage{fancyhdr}
\usepackage{titlesec}
\usepackage{xcolor}
\usepackage{needspace}
\usepackage{xurl}
\usepackage{enumitem}
\setlist[itemize]{leftmargin=2em,itemsep=3pt,topsep=3pt}
\setlength{\parindent}{1.6em}
\setlength{\parskip}{0.25em}
\raggedbottom
\widowpenalty=10000
\clubpenalty=10000
\displaywidowpenalty=10000
\emergencystretch=3em
\definecolor{primary}{HTML}{2B6259}
\definecolor{muted}{HTML}{546671}
\definecolor{sand}{HTML}{F3F6F5}
\titleformat{\section}{\color{primary}\normalfont\bfseries\Large}{\thesection}{1em}{}
\titleformat{\subsection}{\color{primary}\normalfont\bfseries\large}{\thesubsection}{1em}{}
\titleformat{\subsubsection}{\color{primary}\normalfont\bfseries\normalsize}{\thesubsubsection}{1em}{}
\titlespacing*{\section}{0pt}{2.6ex}{1.2ex}
\titlespacing*{\subsection}{0pt}{2.0ex}{0.8ex}
\titlespacing*{\subsubsection}{0pt}{1.4ex}{0.5ex}
\pagestyle{fancy}
\fancyhf{}
\fancyhead[L]{\small\color{muted} 考研英语（一）\quad 学生讲义}
\fancyhead[R]{\small\color{muted} 语言 · 阅读 · 命题 · 表达}
\fancyfoot[C]{\small\thepage}
\renewcommand{\headrulewidth}{0.35pt}
\renewcommand{\footrulewidth}{0pt}
\setlength{\headheight}{16pt}
\renewenvironment{quote}{\begin{tcolorbox}[enhanced,breakable,colback=sand,colframe=primary!22,boxrule=0.4pt,arc=2pt,left=2.6mm,right=2.6mm,top=1.6mm,bottom=1.6mm]\itshape}{\end{tcolorbox}}
\ifLuaTeX
  \usepackage{selnolig}  % disable illegal ligatures
\fi
\usepackage{bookmark}
\IfFileExists{xurl.sty}{\usepackage{xurl}}{} % add URL line breaks if available
\urlstyle{same}
\hypersetup{
  pdftitle={考研英语（一） · L003},
  pdfauthor={英语能力 · 文章理解 · 考题证据 · 迁移表达},
  hidelinks,
  pdfcreator={LaTeX via pandoc}}

\title{考研英语（一） · L003}
\usepackage{etoolbox}
\makeatletter
\providecommand{\subtitle}[1]{% add subtitle to \maketitle
  \apptocmd{\@title}{\par {\large #1 \par}}{}{}
}
\makeatother
\subtitle{工作与职业｜学历价值与劳动保障}
\author{英语能力 · 文章理解 · 考题证据 · 迁移表达}
\date{}

\begin{document}
\maketitle

{
\setcounter{tocdepth}{2}
\tableofcontents
}
\textbf{读懂评价与制度，而不是把作者的主张当成事实}

英语一中的职场文章不仅考察行业词汇。作者常从一项被普遍接受的观念写起，接着质疑它的适用范围，再使用数据、例子或另一国的制度说明自己的判断。

本册的两篇真题分别讨论``学位能为未来提供多大保障''和``劳动保护适用于不同岗位时可能出现什么后果''。前者是观点评论，后者包含经济政策立场。我们既要看清英文表达，也要辨认：哪些是事实、作者怎样推论、题目又允许我们作出多大推断。

\subsubsection{本册学完以后，你应该能够}\label{ux672cux518cux5b66ux5b8cux4ee5ux540eux4f60ux5e94ux8be5ux80fdux591f}

\begin{itemize}
\item
  准确理解抽象评价（be found wanting /
  devalue）与资格、能力、岗位之间的关系。
\item
  学会读懂多重 that
  从句、否定转向、适用条件、反常结果与数据在论证中的作用。
\item
  对试卷中的经济判断保持``作者主张---证据---可推出范围''之间的区分。
\item
  能够用英文说明职业路径的差异，撰写兼顾权利与效率的短段落，而不是复述题目。
\end{itemize}

\textbf{使用顺序}：先做导入判断；阅读外部真题并独立答题；使用本册的原文回访任务和语法词汇解释；完成原创阅读与专项训练；最后看答案，再看中英对照附录。完整原卷不在此重印。

\begin{quote}
\textbf{材料身份}　本册出现的外部试题为指定年份的考研英语（一）真题，参考答案在文末注明为``核对后的非官方参考答案''。所有原创例句、迁移阅读和练习均为本册自编，不属于历年真题。
\end{quote}

\section{进入阅读之前｜你能识别这些区别吗}\label{ux8fdbux5165ux9605ux8bfbux4e4bux524dux4f60ux80fdux8bc6ux522bux8fd9ux4e9bux533aux522bux5417}

先读四句英文，不查中文翻译。圈出最可能影响判断的词组，回答后面的中文问题。对照译文统一放在最后。

\textbf{诊断 01}

\begin{quote}
A qualification can open a door without guaranteeing a career.
\end{quote}

说清 open a door 与 guarantee a career 断言强度的差别。

\begin{center}\rule{0.5\linewidth}{0.5pt}\end{center}

\textbf{诊断 02}

\begin{quote}
The promise was found wanting when graduates entered the job market.
\end{quote}

found wanting 中的 wanting 是动词``想要''吗？

\begin{center}\rule{0.5\linewidth}{0.5pt}\end{center}

\textbf{诊断 03}

\begin{quote}
The policy was intended to protect employees, yet its costs fell
unevenly.
\end{quote}

区分制度初衷和制度效果，并说明yet的作用。

\begin{center}\rule{0.5\linewidth}{0.5pt}\end{center}

\textbf{诊断 04}

\begin{quote}
To the extent that the rule delays hiring, new applicants may be
affected.
\end{quote}

to the extent that 是``全部情况''还是限定成立的条件？

\begin{center}\rule{0.5\linewidth}{0.5pt}\end{center}

如果读不顺，先区分障碍：（1）词义不明 （2）句子关系不清
（3）知道句意但不知道为什么作者这样说
（4）文章懂了却不确定选项。下面四个单元会分别处理。

\section{单元一｜大学学位：一种承诺何时不再充分}\label{ux5355ux5143ux4e00ux5927ux5b66ux5b66ux4f4dux4e00ux79cdux627fux8bfaux4f55ux65f6ux4e0dux518dux5145ux5206}

\textbf{真实试卷}：2022 年英语（一） Reading Text
2（26---30题）。请先打开\href{https://ac-wiki.org/general-skills/study/assets/2022-201.pdf}{可查阅的试卷版本}，标注段号、题号和自己初次选择的答案。

先阅读，不立即给学位是否``值得''的个人观点。记录作者真正质疑的是哪一种承诺；寻找原文明确承认学位仍有价值的位置。

\subsubsection{第一次阅读｜留给自己的证据笔记}\label{ux7b2cux4e00ux6b21ux9605ux8bfbux7559ux7ed9ux81eaux5df1ux7684ux8bc1ux636eux7b14ux8bb0}

请先尝试在纸上完成：

\begin{itemize}
\tightlist
\item
  不看选项，用一句中文写出本文要讨论的问题。
\item
  哪一句/哪一段把``描述对象''转向``提出判断''？
\item
  对最没把握的一题，先在英文原文画出证据，写出自己的英文简答，然后才比较四个选项。
\end{itemize}

\subsubsection{段落阅读实验｜先观察，再返回原文}\label{ux6bb5ux843dux9605ux8bfbux5b9eux9a8cux5148ux89c2ux5bdfux518dux8fd4ux56deux539fux6587}

下列句子是\textbf{原创平行句}，不是从真题抄来的；它们用来演示理解方法。每组请先解释平行句，再回到指定试卷段落完成任务。题目反馈集中在答案部分。

\paragraph{D1 ·
先确立问题，不急着替作者作结论}\label{d1-ux5148ux786eux7acbux95eeux9898ux4e0dux6025ux7740ux66ffux4f5cux8005ux4f5cux7ed3ux8bba}

首段写新一代申请大学时应重新衡量学位意义。consider不是已经决定，weigh
up把选择置于比较之中。

\textbf{原创平行句}

\begin{quote}
Before choosing a course, applicants should weigh up both costs and
possible benefits.
\end{quote}

\textbf{返回原文的任务}：第一段要读者重新衡量什么？能否推出作者主张所有年轻人都不要上大学？

\textbf{自己留下的一句话答案}：

\rule{0.75\linewidth}{0.35pt}

\paragraph{D2 ·
评价不是``想要''：理解承诺被检验}\label{d2-ux8bc4ux4ef7ux4e0dux662fux60f3ux8981ux7406ux89e3ux627fux8bfaux88abux68c0ux9a8c}

第二段的that promise
回指一种稳定的社会预期。后接学位普及、相对贬值和就业结构的统计，构成``观点---说明---证据''。

\textbf{原创平行句}

\begin{quote}
The training was useful, but the promise of automatic promotion was
found wanting.
\end{quote}

\textbf{返回原文的任务}：找出that
promise的内容。统计数字是作者的最终结论还是支撑判断的证据？

\textbf{自己留下的一句话答案}：

\rule{0.75\linewidth}{0.35pt}

\paragraph{D3 ·
让步的精度：不是无价值，而非人人适合}\label{d3-ux8ba9ux6b65ux7684ux7cbeux5ea6ux4e0dux662fux65e0ux4ef7ux503cux800cux975eux4ebaux4ebaux9002ux5408}

第三段明确保留取得学位的意义，同时强调其并不适合每个人，存在其他路径。not\ldots but
rather不能削成简单否定。

\textbf{原创平行句}

\begin{quote}
This is not to deny the value of university, but to question whether one
route suits everyone.
\end{quote}

\textbf{返回原文的任务}：找出P3否定的误读与作者真正肯定的判断。

\textbf{自己留下的一句话答案}：

\rule{0.75\linewidth}{0.35pt}

\paragraph{D4 ·
雇主标准的变化与证据边界}\label{d4-ux96c7ux4e3bux6807ux51c6ux7684ux53d8ux5316ux4e0eux8bc1ux636eux8fb9ux754c}

第四段谈雇主对部分非毕业生人员的评价及取消个别岗位学位要求。某些雇主做法不能扩张为全部市场。

\textbf{原创平行句}

\begin{quote}
Some employers removed degree requirements for roles that relied on
practical skills.
\end{quote}

\textbf{返回原文的任务}：P4中的some/certain如何限制答案？取消学位门槛是否等于取消能力考核？

\textbf{自己留下的一句话答案}：

\rule{0.75\linewidth}{0.35pt}

\paragraph{D5 ·
专门知识与一般能力并非二选一}\label{d5-ux4e13ux95e8ux77e5ux8bc6ux4e0eux4e00ux822cux80fdux529bux5e76ux975eux4e8cux9009ux4e00}

第五段转向仍选择高等教育的人，强调专长可能为他们创造额外价值。作者不是说任何人都应该连读两个学位。

\textbf{原创平行句}

\begin{quote}
A general education can be strengthened by specialised training in a
chosen field.
\end{quote}

\textbf{返回原文的任务}：作者在哪一类人身上强调进一步专业学习？这个建议是否适用于所有人？

\textbf{自己留下的一句话答案}：

\rule{0.75\linewidth}{0.35pt}

\paragraph{D6 ·
职业轨迹是持续变化的，不是单一身份}\label{d6-ux804cux4e1aux8f68ux8ff9ux662fux6301ux7eedux53d8ux5316ux7684ux4e0dux662fux5355ux4e00ux8eabux4efd}

最后两段强调持续提升技能与多次转换岗位的可能，并对老一代的``学位即身份''思维作对照。就业数字是文章使用的估计，不是个人必然会经历的次数。

\textbf{原创平行句}

\begin{quote}
A degree may begin a career trajectory without defining the whole
journey.
\end{quote}

\textbf{返回原文的任务}：最后一句对``我就是某某专业的人''的表达持何种态度？

\textbf{自己留下的一句话答案}：

\rule{0.75\linewidth}{0.35pt}

\subsubsection{词汇与表达｜从具体误解入手}\label{ux8bcdux6c47ux4e0eux8868ux8fbeux4eceux5177ux4f53ux8befux89e3ux5165ux624b}

这不是统一格式的``词义抄写表''。请先理解词在本文中的意义，关注它改变了哪一层断言，再试着不看中文译出原创例句。无需把所有义项一口气背完。

\paragraph{01　weigh up}\label{weigh-up}

\textbf{本篇核心用法}　phr.
权衡；认真比较选项。　\textbf{典型组合}　weigh up options / costs and
benefits。

\textbf{容易误解的地方}　强调尚在决策过程中，非``称重''字面义。

\textbf{原创语境句}

\begin{quote}
Applicants must weigh up tuition costs against opportunities for
practical training.
\end{quote}

\textbf{再想一步}　比较权衡是否已经说明最后要选择哪项？

\paragraph{02　be found wanting}\label{be-found-wanting}

\textbf{本篇核心用法}　phr.
被证明不足；未达到预期。　\textbf{典型组合}　a promise was found
wanting。

\textbf{容易误解的地方}　found是find的过去分词；wanting作表语指``不足''，不是``正在想要''。

\textbf{原创语境句}

\begin{quote}
The claim that one certificate guaranteed a job was eventually found
wanting.
\end{quote}

\textbf{再想一步}　应当评价的是哪一个``承诺/说法''？

\paragraph{03　devalue}\label{devalue}

\textbf{本篇核心用法}　v.
使价值下降；降低相对重要性。　\textbf{典型组合}　devalue a qualification
/ currency。

\textbf{容易误解的地方}　在考试里指普及导致的区分度降低，不等于其使用价值归零。

\textbf{原创语境句}

\begin{quote}
When nearly everyone held the badge, it no longer distinguished
applicants and became devalued.
\end{quote}

\textbf{再想一步}　区分度下降和完全无用的差别是什么？

\paragraph{04　social mobility}\label{social-mobility}

\textbf{本篇核心用法}　n.~社会流动；社会阶层位置变化。　\textbf{典型组合}　a
route to social mobility。

\textbf{容易误解的地方}　指地位或机会向上流动，不是交通工具的``移动''。

\textbf{原创语境句}

\begin{quote}
Affordable training can offer a route to social mobility for some
families.
\end{quote}

\textbf{再想一步}　some为何是重要限定？

\paragraph{05　inevitable}\label{inevitable}

\textbf{本篇核心用法}　adj. 不可避免的。　\textbf{典型组合}　an
inevitable outcome / not inevitable。

\textbf{容易误解的地方}　以必然性表述；not inevitable不等于never
happens。

\textbf{原创语境句}

\begin{quote}
The move from secondary school to university is not inevitable for every
learner.
\end{quote}

\textbf{再想一步}　not inevitable 与 impossible 是否同义？

\paragraph{06　predecessor}\label{predecessor}

\textbf{本篇核心用法}　n.~前辈；前任。　\textbf{典型组合}　learn from
predecessors。

\textbf{容易误解的地方}　表示时间或职位上的前一群体，不能与同龄人混淆。

\textbf{原创语境句}

\begin{quote}
Young engineers often learn from the mistakes of their predecessors.
\end{quote}

\textbf{再想一步}　例句是否断言前辈总是正确？

\paragraph{07　committed}\label{committed}

\textbf{本篇核心用法}　adj.
忠诚、投入的；致力于。　\textbf{典型组合}　committed to a role /
committed employees。

\textbf{容易误解的地方}　描述投入态度；不能仅从学历推断员工是否投入。

\textbf{原创语境句}

\begin{quote}
A committed trainee can make progress without following a conventional
degree route.
\end{quote}

\textbf{再想一步}　投入程度是否由学位有无直接决定？

\paragraph{08　scrap}\label{scrap}

\textbf{本篇核心用法}　v. 废止；取消。　\textbf{典型组合}　scrap a
requirement / plan。

\textbf{容易误解的地方}　在P4是取消某些岗位学位要求，而非废除全部资格标准。

\textbf{原创语境句}

\begin{quote}
The employer scrapped a degree requirement but retained the skills test.
\end{quote}

\textbf{再想一步}　为何 but 是理解改革范围的关键？

\paragraph{09　generalist}\label{generalist}

\textbf{本篇核心用法}　n.~通才。　\textbf{典型组合}　a generalist / a
specialist。

\textbf{容易误解的地方}　可与specialist形成教学对照；并不等于什么都不会。

\textbf{原创语境句}

\begin{quote}
A generalist may connect ideas across departments but still need deep
expertise in one task.
\end{quote}

\textbf{再想一步}　通才与专门知识一定冲突吗？

\paragraph{10　postgraduate}\label{postgraduate}

\textbf{本篇核心用法}　n./adj.
研究生；研究生阶段的。　\textbf{典型组合}　postgraduate study /
postgraduate students。

\textbf{容易误解的地方}　英式用法常作研究生阶段；不能直接等同就业保证。

\textbf{原创语境句}

\begin{quote}
Postgraduate study may deepen expertise, although its value varies by
field.
\end{quote}

\textbf{再想一步}　后半句限制了哪一项过强推断？

\paragraph{11　upskill}\label{upskill}

\textbf{本篇核心用法}　v.
提升技能；进修新技能。　\textbf{典型组合}　upskill throughout a career。

\textbf{容易误解的地方}　强调职业持续学习，不单指学历升级。

\textbf{原创语境句}

\begin{quote}
Workers can upskill through short courses while remaining in employment.
\end{quote}

\textbf{再想一步}　这是否要求必须脱产取得新学位？

\paragraph{12　employable}\label{employable}

\textbf{本篇核心用法}　adj.
具备受雇条件的；有就业竞争力的。　\textbf{典型组合}　remain employable /
employability。

\textbf{容易误解的地方}　ability to get work 不等于 guarantee of getting
one particular job。

\textbf{原创语境句}

\begin{quote}
Technical knowledge and communication skills can help graduates remain
employable.
\end{quote}

\textbf{再想一步}　有能力被雇用等于获得聘用承诺吗？

\paragraph{13　trajectory}\label{trajectory}

\textbf{本篇核心用法}　n.~发展轨迹；路径。　\textbf{典型组合}　career
trajectory / growth trajectory。

\textbf{容易误解的地方}　强调随时间展开的变化路线，不是固定岗位。

\textbf{原创语境句}

\begin{quote}
Her career trajectory included teaching, design and later project
management.
\end{quote}

\textbf{再想一步}　为何不能用第一份工作定义全部career trajectory？

\paragraph{14　fulfillment}\label{fulfillment}

\textbf{本篇核心用法}　n.~满足感；成就感；实现。　\textbf{典型组合}　personal
fulfillment / fulfillment of a promise。

\textbf{容易误解的地方}　真题取``个人满足感''，其他语境可指履行承诺。

\textbf{原创语境句}

\begin{quote}
Some people value personal fulfillment more than a prestigious job
title.
\end{quote}

\textbf{再想一步}　它与salary是同一种评价维度吗？

\subsubsection{与上一册连接的旧词}\label{ux4e0eux4e0aux4e00ux518cux8fdeux63a5ux7684ux65e7ux8bcd}

\begin{itemize}
\item
  qualified（L001）---符合条件的；与qualification名词关联，但不强行等同。
\item
  transition（L001）---转变过程；关联从校园到岗位的过渡。
\end{itemize}

\subsubsection{回到试卷：给初次判断一次修正机会}\label{ux56deux5230ux8bd5ux5377ux7ed9ux521dux6b21ux5224ux65adux4e00ux6b21ux4feeux6b63ux673aux4f1a}

再做一遍这一篇真题的五道题；不要直接照抄第一次答案。对每道题写出：\textbf{原文哪一处提供证据？题干要求概括事实、判断态度还是进行推断？最诱人的错误项改变了什么？}

\begin{itemize}
\item
  \textbf{第26题}　证据段：\_\_\_\_　简答另写：\_\_\_\_　最易误选项：\_\_\_\_
\item
  \textbf{第27题}　证据段：\_\_\_\_　简答另写：\_\_\_\_　最易误选项：\_\_\_\_
\item
  \textbf{第28题}　证据段：\_\_\_\_　简答另写：\_\_\_\_　最易误选项：\_\_\_\_
\item
  \textbf{第29题}　证据段：\_\_\_\_　简答另写：\_\_\_\_　最易误选项：\_\_\_\_
\item
  \textbf{第30题}　证据段：\_\_\_\_　简答另写：\_\_\_\_　最易误选项：\_\_\_\_
\end{itemize}

\section{单元二｜劳动保障：目的与后果之间有几层推理}\label{ux5355ux5143ux4e8cux52b3ux52a8ux4fddux969cux76eeux7684ux4e0eux540eux679cux4e4bux95f4ux6709ux51e0ux5c42ux63a8ux7406}

\textbf{真实试卷}：2022 年英语（一） Reading Text
4（36---40题）。请先打开\href{https://ac-wiki.org/general-skills/study/assets/2022-201.pdf}{可查阅的试卷版本}，标注段号、题号和自己初次选择的答案。

第一次阅读暂时不评价政策对错。画三栏：法律意图、作者声称的经济机制、文章引用的依据。把每个判断放到正确栏中。

\subsubsection{第一次阅读｜留给自己的证据笔记}\label{ux7b2cux4e00ux6b21ux9605ux8bfbux7559ux7ed9ux81eaux5df1ux7684ux8bc1ux636eux7b14ux8bb0-1}

请先尝试在纸上完成：

\begin{itemize}
\tightlist
\item
  不看选项，用一句中文写出本文要讨论的问题。
\item
  哪一句/哪一段把``描述对象''转向``提出判断''？
\item
  对最没把握的一题，先在英文原文画出证据，写出自己的英文简答，然后才比较四个选项。
\end{itemize}

\subsubsection{段落阅读实验｜先观察，再返回原文}\label{ux6bb5ux843dux9605ux8bfbux5b9eux9a8cux5148ux89c2ux5bdfux518dux8fd4ux56deux539fux6587-1}

下列句子是\textbf{原创平行句}，不是从真题抄来的；它们用来演示理解方法。每组请先解释平行句，再回到指定试卷段落完成任务。题目反馈集中在答案部分。

\paragraph{E1 ·
规则保护谁：先读法律的原初目的}\label{e1-ux89c4ux5219ux4fddux62a4ux8c01ux5148ux8bfbux6cd5ux5f8bux7684ux539fux521dux76eeux7684}

第一、二段描述法规限制无正当理由的解雇，并要求程序公平。题目第36问的是制度设计目的，不是宏观经济最终影响。

\textbf{原创平行句}

\begin{quote}
A dismissal rule may require an employer to explain the reason and
follow a fair procedure.
\end{quote}

\textbf{返回原文的任务}：原文区别``必须给出理由''与``必须遵守程序''了吗？

\textbf{自己留下的一句话答案}：

\rule{0.75\linewidth}{0.35pt}

\paragraph{E2 ·
讨论适用对象，而非否定所有保护}\label{e2-ux8ba8ux8bbaux9002ux7528ux5bf9ux8c61ux800cux975eux5426ux5b9aux6240ux6709ux4fddux62a4}

第三段作者开始质疑把同样的保护程序套用于高薪管理者。必须看见限定语，而不能改写成``劳动保护必然妨碍发展''。

\textbf{原创平行句}

\begin{quote}
A rule designed for new workers may have different effects when applied
to senior executives.
\end{quote}

\textbf{返回原文的任务}：P3的条件限定涉及哪些岗位/收入群体？

\textbf{自己留下的一句话答案}：

\rule{0.75\linewidth}{0.35pt}

\paragraph{E3 ·
意外结果为何需要中间环节}\label{e3-ux610fux5916ux7ed3ux679cux4e3aux4f55ux9700ux8981ux4e2dux95f4ux73afux8282}

第四段作者提出反常观点：保护工作岗位的法律可能带来某些工作风险。但结论依赖管理效率与企业表现等中介机制。

\textbf{原创平行句}

\begin{quote}
A policy can be well intended yet produce costs through the choices it
changes.
\end{quote}

\textbf{返回原文的任务}：怎样补全``保护限制解雇 -\textgreater{}
管理能力调整受阻 -\textgreater{} 企业表现影响''的因果链？

\textbf{自己留下的一句话答案}：

\rule{0.75\linewidth}{0.35pt}

\paragraph{E4 · 注意 to the extent
的证据边界}\label{e4-ux6ce8ux610f-to-the-extent-ux7684ux8bc1ux636eux8fb9ux754c}

第五段中to the
extent引导条件强度：只有在保护法确实限制某类管理调整的范围内，作者才推导工资和效率后果。

\textbf{原创平行句}

\begin{quote}
To the extent that hiring becomes more costly, a firm may recruit more
cautiously.
\end{quote}

\textbf{返回原文的任务}：这句是否断言每家公司都减少全部员工的工资？

\textbf{自己留下的一句话答案}：

\rule{0.75\linewidth}{0.35pt}

\paragraph{E5 ·
从论据到推论：不可把引用当成万能证明}\label{e5-ux4eceux8bbaux636eux5230ux63a8ux8bbaux4e0dux53efux628aux5f15ux7528ux5f53ux6210ux4e07ux80fdux8bc1ux660e}

文章引用生产率委员会关于管理能力的判断，为``低水平管理影响生产率''提供材料，但它并未做完所有经济因果的独立实验。

\textbf{原创平行句}

\begin{quote}
The report identified weak management as one possible source of low
productivity.
\end{quote}

\textbf{返回原文的任务}：被引用资料直接支持什么，哪些环节仍由作者推进？

\textbf{自己留下的一句话答案}：

\rule{0.75\linewidth}{0.35pt}

\paragraph{E6 ·
跨国对照是方案，不是普遍事实}\label{e6-ux8de8ux56fdux5bf9ux7167ux662fux65b9ux6848ux4e0dux662fux666eux904dux4e8bux5b9e}

末段提到澳大利亚高收入门槛例外，新西兰相关提案未通过。作者使用制度对照推论企业负担可能下降，不是证明该安排适合所有国家。

\textbf{原创平行句}

\begin{quote}
An income threshold excludes some senior staff from a rule while leaving
others covered.
\end{quote}

\textbf{返回原文的任务}：如何理解threshold在这里的作用？

\textbf{自己留下的一句话答案}：

\rule{0.75\linewidth}{0.35pt}

\subsubsection{词汇与表达｜从具体误解入手}\label{ux8bcdux6c47ux4e0eux8868ux8fbeux4eceux5177ux4f53ux8befux89e3ux5165ux624b-1}

这不是统一格式的``词义抄写表''。请先理解词在本文中的意义，关注它改变了哪一层断言，再试着不看中文译出原创例句。无需把所有义项一口气背完。

\paragraph{01　grievance}\label{grievance}

\textbf{本篇核心用法}　n.~申诉；不满；冤屈。　\textbf{典型组合}　personal
grievance / file a grievance。

\textbf{容易误解的地方}　在劳动法语境常为雇员提出的正式申诉程序。

\textbf{原创语境句}

\begin{quote}
The employee filed a grievance after receiving no explanation for the
dismissal.
\end{quote}

\textbf{再想一步}　这只是私人抱怨，还是可能启动正式程序？

\paragraph{02　provision}\label{provision}

\textbf{本篇核心用法}　n.~条款；规定；供应。　\textbf{典型组合}　legal
provisions / provisions of an act。

\textbf{容易误解的地方}　本篇是法律条款，不是物资供给；按语境判断。

\textbf{原创语境句}

\begin{quote}
The provision requires employers to keep a written record of the
decision.
\end{quote}

\textbf{再想一步}　它指机构还是成文规定？

\paragraph{03　dismissal}\label{dismissal}

\textbf{本篇核心用法}　n.~解雇；驳回。　\textbf{典型组合}　unjustified
dismissal / unfair dismissal。

\textbf{容易误解的地方}　劳动关系语境是终止雇佣，不是``否定某项观点''。

\textbf{原创语境句}

\begin{quote}
A dismissal can be lawful only if the required procedures are followed.
\end{quote}

\textbf{再想一步}　can 与 only if 对结论加了什么限制？

\paragraph{04　justify}\label{justify}

\textbf{本篇核心用法}　v.
证明\ldots\ldots 正当；给出充分理由。　\textbf{典型组合}　justify a
decision / be justified。

\textbf{容易误解的地方}　理由存在与程序合规是不同要求。

\textbf{原创语境句}

\begin{quote}
The manager could not justify the decision with reliable evidence.
\end{quote}

\textbf{再想一步}　有理由是否自动满足全部程序要求？

\paragraph{05　procedural}\label{procedural}

\textbf{本篇核心用法}　adj. 程序上的。　\textbf{典型组合}　procedural
fairness / procedural rules。

\textbf{容易误解的地方}　真题出现副词procedurally；词性转换不改变``程序''概念。

\textbf{原创语境句}

\begin{quote}
Procedural fairness requires both notice and a chance to respond.
\end{quote}

\textbf{再想一步}　这里讨论结果好坏还是实施过程公正？

\paragraph{06　safeguard}\label{safeguard}

\textbf{本篇核心用法}　n./v. 保障措施；保护。　\textbf{典型组合}　legal
safeguard / safeguard rights。

\textbf{容易误解的地方}　不同于 guarantee
全面保证；是降低风险的制度安排。

\textbf{原创语境句}

\begin{quote}
A transparent review process can safeguard employees against arbitrary
decisions.
\end{quote}

\textbf{再想一步}　可以推出所有不公平决定将消失吗？

\paragraph{07　arbitrary}\label{arbitrary}

\textbf{本篇核心用法}　adj.
任意的；缺乏充分理由的。　\textbf{典型组合}　arbitrary dismissal /
arbitrary conduct。

\textbf{容易误解的地方}　不是``随机抽样''的中性统计义，而是负面评价。

\textbf{原创语境句}

\begin{quote}
The new procedure aims to prevent arbitrary changes to pay.
\end{quote}

\textbf{再想一步}　``旨在预防''与``已经杜绝''相同吗？

\paragraph{08　constrain}\label{constrain}

\textbf{本篇核心用法}　v. 限制；约束。　\textbf{典型组合}　constrain
firms / constrained by costs。

\textbf{容易误解的地方}　约束可能改变可选行为，不必然决定最终结果。

\textbf{原创语境句}

\begin{quote}
High training costs may constrain the number of apprentices a firm can
hire.
\end{quote}

\textbf{再想一步}　may为何不能直接改为will always？

\paragraph{09　productivity}\label{productivity}

\textbf{本篇核心用法}　n.~生产率；单位投入产出。　\textbf{典型组合}　labour
productivity / boost productivity。

\textbf{容易误解的地方}　与利润、工资不同；它们可能相关但不是同义。

\textbf{原创语境句}

\begin{quote}
The factory improved productivity without increasing the number of
working hours.
\end{quote}

\textbf{再想一步}　为什么不能直接译成``总产量增加''？

\paragraph{10　paradoxically}\label{paradoxically}

\textbf{本篇核心用法}　adv.
看似矛盾地；出人意料地。　\textbf{典型组合}　paradoxically, \ldots。

\textbf{容易误解的地方}　指出结果与政策初衷之间反直觉的关系，但不自带科学验证。

\textbf{原创语境句}

\begin{quote}
Paradoxically, a rule meant to speed up hiring created more paperwork.
\end{quote}

\textbf{再想一步}　哪一个是初衷，哪一个是结果？

\paragraph{11　marginal}\label{marginal}

\textbf{本篇核心用法}　adj.
边缘的；差一点达到条件的。　\textbf{典型组合}　marginal candidate /
marginal change。

\textbf{容易误解的地方}　经济语境可指较难竞争/位于边缘的位置，不能机械等于``微不足道''。

\textbf{原创语境句}

\begin{quote}
A marginal applicant may be excluded when recruitment costs rise.
\end{quote}

\textbf{再想一步}　被排除一定是因为能力为零吗？

\paragraph{12　stringent}\label{stringent}

\textbf{本篇核心用法}　adj.
严格的；要求苛刻的。　\textbf{典型组合}　stringent rules / stringent
standards。

\textbf{容易误解的地方}　描述规则强度；是否合理须另由证据评价。

\textbf{原创语境句}

\begin{quote}
Stringent reporting rules may improve transparency but raise
administrative costs.
\end{quote}

\textbf{再想一步}　严格规则是否总是坏事？

\paragraph{13　threshold}\label{threshold}

\textbf{本篇核心用法}　n.~门槛；阈值。　\textbf{典型组合}　income
threshold / eligibility threshold。

\textbf{容易误解的地方}　用于划分适用范围；与minimum
wage的含义完全不同。

\textbf{原创语境句}

\begin{quote}
Employees above the income threshold are covered by a different review
procedure.
\end{quote}

\textbf{再想一步}　收入门槛决定覆盖范围，而非决定每人应拿多少工资。

\paragraph{14　unwieldy}\label{unwieldy}

\textbf{本篇核心用法}　adj.
笨重难用的；程序繁琐的。　\textbf{典型组合}　an unwieldy procedure /
mechanism。

\textbf{容易误解的地方}　法律机制抽象义为复杂且难实施，不是物理物体巨大。

\textbf{原创语境句}

\begin{quote}
The proposed appeal system was too unwieldy for small employers to use.
\end{quote}

\textbf{再想一步}　繁琐是否与不必要同义？

\subsubsection{与上一册连接的旧词}\label{ux4e0eux4e0aux4e00ux518cux8fdeux63a5ux7684ux65e7ux8bcd-1}

\begin{itemize}
\item
  incentive（L001）---行为激励与政策作用机制。
\item
  regulate（S001）---规制某种行为，注意与protect的区别。
\item
  consequence（S001）---后果；单凭词面不说明因果已经得到实验验证。
\end{itemize}

\subsubsection{回到试卷：给初次判断一次修正机会}\label{ux56deux5230ux8bd5ux5377ux7ed9ux521dux6b21ux5224ux65adux4e00ux6b21ux4feeux6b63ux673aux4f1a-1}

再做一遍这一篇真题的五道题；不要直接照抄第一次答案。对每道题写出：\textbf{原文哪一处提供证据？题干要求概括事实、判断态度还是进行推断？最诱人的错误项改变了什么？}

\begin{itemize}
\item
  \textbf{第36题}　证据段：\_\_\_\_　简答另写：\_\_\_\_　最易误选项：\_\_\_\_
\item
  \textbf{第37题}　证据段：\_\_\_\_　简答另写：\_\_\_\_　最易误选项：\_\_\_\_
\item
  \textbf{第38题}　证据段：\_\_\_\_　简答另写：\_\_\_\_　最易误选项：\_\_\_\_
\item
  \textbf{第39题}　证据段：\_\_\_\_　简答另写：\_\_\_\_　最易误选项：\_\_\_\_
\item
  \textbf{第40题}　证据段：\_\_\_\_　简答另写：\_\_\_\_　最易误选项：\_\_\_\_
\end{itemize}

\section{句子工坊｜精准理解比翻译顺口更重要}\label{ux53e5ux5b50ux5de5ux574aux7cbeux51c6ux7406ux89e3ux6bd4ux7ffbux8bd1ux987aux53e3ux66f4ux91cdux8981}

下面六句均为原创，请先划分主干，再圈出限定词、连接词或修饰范围。\textbf{先用英语或中文说清命题，再翻译}；不要直接猜中文意思。

\subsubsection{S1｜句子观察}\label{s1ux53e5ux5b50ux89c2ux5bdf}

\begin{quote}
The promise that a degree would secure a job was questioned as more
graduates entered the labour market.
\end{quote}

\textbf{观察任务}　找出that从句解释哪一个名词；as是时间伴随还是原因必须结合语境？

\textbf{独立翻译}　

\rule{0.8\linewidth}{0.35pt}

\subsubsection{S2｜句子观察}\label{s2ux53e5ux5b50ux89c2ux5bdf}

\begin{quote}
This is not to claim that qualifications are useless, but to ask whether
they are sufficient.
\end{quote}

\textbf{观察任务}　辨认not to\ldots but to\ldots 的立场转向。

\textbf{独立翻译}　

\rule{0.8\linewidth}{0.35pt}

\subsubsection{S3｜句子观察}\label{s3ux53e5ux5b50ux89c2ux5bdf}

\begin{quote}
An employer may remove a degree requirement while still demanding
evidence of practical competence.
\end{quote}

\textbf{观察任务}　while的关系是什么；取消A是否等于取消B？

\textbf{独立翻译}　

\rule{0.8\linewidth}{0.35pt}

\subsubsection{S4｜句子观察}\label{s4ux53e5ux5b50ux89c2ux5bdf}

\begin{quote}
Although the rule was intended to protect workers, its effects varied
with the positions to which it applied.
\end{quote}

\textbf{观察任务}　划出规则目的、应用对象与效果差异。

\textbf{独立翻译}　

\rule{0.8\linewidth}{0.35pt}

\subsubsection{S5｜句子观察}\label{s5ux53e5ux5b50ux89c2ux5bdf}

\begin{quote}
To the extent that the procedure adds costs, firms may become more
cautious about hiring.
\end{quote}

\textbf{观察任务}　这句话含几个收窄结论的语言线索？

\textbf{独立翻译}　

\rule{0.8\linewidth}{0.35pt}

\subsubsection{S6｜句子观察}\label{s6ux53e5ux5b50ux89c2ux5bdf}

\begin{quote}
A report linking management quality to productivity does not by itself
prove that every dismissal rule reduces wages.
\end{quote}

\textbf{观察任务}　区分报告观察与作者外推；说明by itself的作用。

\textbf{独立翻译}　

\rule{0.8\linewidth}{0.35pt}

\textbf{读者提示}：如果你能流畅译出一个句子，却无法说明``谁在什么条件下断言了什么''，仍需要回到连词、限定和先行词。真正重要的是恢复命题关系。

\section{阅读迁移｜两篇没有见过的英文文章}\label{ux9605ux8bfbux8fc1ux79fbux4e24ux7bc7ux6ca1ux6709ux89c1ux8fc7ux7684ux82f1ux6587ux6587ux7ae0}

以下均为\textbf{原创训练文章}。先按真实阅读的方式独立完成五道选择题和两道无选项简答，再读答案。这里的目标不是猜同一个题型，而是检验词义、篇章和证据判断能否迁移到新题材。

\subsection{Reading A · The Hiring Test That Changed the
Queue}\label{reading-a-the-hiring-test-that-changed-the-queue}

原创迁移阅读：企业以能力测试替代一部分学历门槛，考查证据与政策效果的区别。　\textbf{建议首读限时：7---9分钟；不要先翻到附录看中文。}

\textbf{P1}　When a design company removed the degree requirement from
its entry-level vacancies, the change attracted attention. Managers
explained that many successful employees had learned their skills
outside university. The new advertisements asked applicants to complete
a practical task instead of listing their academic qualifications. The
firm described the policy as a way to widen access, but it offered no
guarantee that everyone would have an equal chance of being hired.

\textbf{P2}　Applications rose sharply during the first month. Some
observers took this as proof that the scheme had succeeded. Yet an
increase in applications was an early response, not evidence about the
quality of recruitment. The company still needed to know whether the
tests identified people who could do the job and whether the extra
assessment time was reasonable.

\textbf{P3}　The practical test also created a new obstacle. Applicants
with access to expensive equipment could complete it quickly, while
others had to borrow a computer or find a quiet room. The absence of a
degree requirement had removed one barrier but left others in place.
After receiving this feedback, the firm began offering equipment at
local libraries and extended the submission deadline.

\textbf{P4}　Six months later, the firm reported that applicants from a
wider range of backgrounds had reached the interview stage. It did not
claim that the new approach had increased long-term productivity.
Managers planned to compare retention rates and performance over several
years before drawing a stronger conclusion. The example suggests that
opening an entrance is a useful first step, but it does not tell us
everything about what happens after people enter.

\subsubsection{Questions ·
独立答题}\label{questions-ux72ecux7acbux7b54ux9898}

\textbf{A1. What replaced the degree requirement?}

\begin{itemize}
\item
  \textbf{A.} A practical task.
\item
  \textbf{B.} A compulsory second degree.
\item
  \textbf{C.} An unpaid two-year contract.
\item
  \textbf{D.} An automatic interview.
\end{itemize}

\textbf{A2. Why was the rise in applications insufficient evidence of
success?}

\begin{itemize}
\item
  \textbf{A.} It occurred several years later.
\item
  \textbf{B.} It said little about recruitment quality.
\item
  \textbf{C.} It was limited to university graduates.
\item
  \textbf{D.} It proved the test unfair.
\end{itemize}

\textbf{A3. Some applicants were disadvantaged because}

\begin{itemize}
\item
  \textbf{A.} the company refused all applications.
\item
  \textbf{B.} the test was about economic theory.
\item
  \textbf{C.} their access to equipment was limited.
\item
  \textbf{D.} libraries had been permanently closed.
\end{itemize}

\textbf{A4. The firm changed its procedure by}

\begin{itemize}
\item
  \textbf{A.} removing the practical test entirely.
\item
  \textbf{B.} reintroducing the degree rule.
\item
  \textbf{C.} promising every applicant a job.
\item
  \textbf{D.} offering equipment and more time.
\end{itemize}

\textbf{A5. The author's central point is that}

\begin{itemize}
\item
  \textbf{A.} wider access is only one part of evaluating reform.
\item
  \textbf{B.} all academic qualifications are worthless.
\item
  \textbf{C.} practical tests always reduce fairness.
\item
  \textbf{D.} early application numbers prove productivity gains.
\end{itemize}

\subsubsection{No-option questions ·
不依赖选项的理解}\label{no-option-questions-ux4e0dux4f9dux8d56ux9009ux9879ux7684ux7406ux89e3}

\textbf{AO1.} Why did the firm avoid claiming higher productivity?

回答时至少指出原文依据及必要限定。

\textbf{AO2.} What trade-off did the practical test introduce?

回答时至少指出原文依据及必要限定。

\subsubsection{读后复盘}\label{ux8bfbux540eux590dux76d8}

闭卷用约40字写出：作者起初介绍的问题、最重要的转向、以及末段的有限结论。然后回到文中画出支撑这三点的句子。

\subsection{Reading B · What a Fair Complaint Procedure Can and Cannot
Do}\label{reading-b-what-a-fair-complaint-procedure-can-and-cannot-do}

原创迁移阅读：新的内部申诉规则与其实际运行，刻意不预设法律好坏。　\textbf{建议首读限时：7---9分钟；不要先翻到附录看中文。}

\textbf{P1}　A medium-sized logistics firm introduced a new complaint
procedure after several employees said that decisions about shifts were
difficult to challenge. Under the procedure, managers had to explain
changes in writing, and employees could request a review by someone who
had not made the original decision. The policy was meant to reduce
arbitrary treatment, not to prevent the company from changing schedules
whenever necessary.

\textbf{P2}　At first, the number of formal complaints increased. Some
managers viewed this as a sign that the workplace had become less
harmonious. An employee representative offered another interpretation:
workers may simply have become more willing to report problems that
previously went unrecorded. Neither interpretation could be established
from the complaint count alone.

\textbf{P3}　The procedure also required time. Supervisors said that
documenting small changes left them less time for daily operations. The
employee representative argued that clear records could prevent larger
disputes later. The disagreement was partly about when the costs and
benefits would appear. A report that examined only the first month would
capture paperwork, but not necessarily avoided disputes six months
later.

\textbf{P4}　The firm therefore agreed to review the policy after a full
year, using several measures: time spent on paperwork, repeat
complaints, staff turnover and employees' confidence in the process. No
single number would settle whether the procedure was fair or efficient.
The more useful question was whether it could protect the right to
challenge a decision without making routine work unnecessarily
difficult.

\subsubsection{Questions ·
独立答题}\label{questions-ux72ecux7acbux7b54ux9898-1}

\textbf{B1. The new procedure was mainly designed to}

\begin{itemize}
\item
  \textbf{A.} abolish all schedule changes.
\item
  \textbf{B.} give workers a way to challenge decisions.
\item
  \textbf{C.} guarantee every complaint would succeed.
\item
  \textbf{D.} replace all supervisors.
\end{itemize}

\textbf{B2. The rise in complaints could mean that}

\begin{itemize}
\item
  \textbf{A.} workers had become more willing to report problems.
\item
  \textbf{B.} all managers were acting illegally.
\item
  \textbf{C.} the procedure had certainly failed.
\item
  \textbf{D.} no workplace problems existed before.
\end{itemize}

\textbf{B3. What did the two sides disagree about in Paragraph 3?}

\begin{itemize}
\item
  \textbf{A.} Whether any record should exist at all.
\item
  \textbf{B.} Who should own the company.
\item
  \textbf{C.} When the procedure's costs and benefits might emerge.
\item
  \textbf{D.} Whether shifts were legal in principle.
\end{itemize}

\textbf{B4. The annual review will measure all except}

\begin{itemize}
\item
  \textbf{A.} staff turnover.
\item
  \textbf{B.} time spent on documentation.
\item
  \textbf{C.} trust in the procedure.
\item
  \textbf{D.} the firm's international market share.
\end{itemize}

\textbf{B5. What is the author's preferred approach?}

\begin{itemize}
\item
  \textbf{A.} Reject it because complaints increased.
\item
  \textbf{B.} Assess both procedural protection and practical costs.
\item
  \textbf{C.} Judge success only by the number of complaints.
\item
  \textbf{D.} Assume all employee reports are inaccurate.
\end{itemize}

\subsubsection{No-option questions ·
不依赖选项的理解}\label{no-option-questions-ux4e0dux4f9dux8d56ux9009ux9879ux7684ux7406ux89e3-1}

\textbf{BO1.} Why can complaint numbers alone mislead?

回答时至少指出原文依据及必要限定。

\textbf{BO2.} How would you distinguish a policy's intention from its
result?

回答时至少指出原文依据及必要限定。

\subsubsection{读后复盘}\label{ux8bfbux540eux590dux76d8-1}

闭卷用约40字写出：作者起初介绍的问题、最重要的转向、以及末段的有限结论。然后回到文中画出支撑这三点的句子。

\section{考试专项衔接｜不只训练四选一阅读}\label{ux8003ux8bd5ux4e13ux9879ux8854ux63a5ux4e0dux53eaux8badux7ec3ux56dbux9009ux4e00ux9605ux8bfb}

以下练习把本册词句带入词汇使用、语篇衔接、翻译和写作。题目均为原创，先完成再看答案。

\subsection{A　语境选词｜结构与词义同时成立}\label{a-ux8bedux5883ux9009ux8bcdux7ed3ux6784ux4e0eux8bcdux4e49ux540cux65f6ux6210ux7acb}

\textbf{词库（无需全部用完）}：weigh up, found wanting, scrap, upskill,
grievance, arbitrary, threshold,
unwieldy。必要时改变动词形态或填入短语的适当形式。

\textbf{V01.} Young applicants should \_\_\_\_\_ the value of formal
study against other options.

\textbf{V02.} The claim of guaranteed employment was \_\_\_\_\_ after
several years of weak results.

\textbf{V03.} The firm decided to \_\_\_\_\_ its degree requirement for
one technical role.

\textbf{V04.} Employees can \_\_\_\_\_ their skills through short
programmes.

\textbf{V05.} The new procedure gives employees a way to file a
\_\_\_\_\_.

\textbf{V06.} A flexible rule can help prevent \_\_\_\_\_ decisions by
management.

\textbf{V07.} The law applies only to employees below a specified income
\_\_\_\_\_.

\textbf{V08.} The proposed approval process proved too \_\_\_\_\_ for
small businesses.

\subsection{B　篇章衔接｜从段落功能而不是重复词作答}\label{b-ux7bc7ux7ae0ux8854ux63a5ux4eceux6bb5ux843dux529fux80fdux800cux4e0dux662fux91cdux590dux8bcdux4f5cux7b54}

\textbf{篇章衔接：从一个招聘数字看两种解释}。根据 G1、G2 前后的信息，从
A---D 选择适当的句子；每项最多使用一次。

A college announced that 82 percent of its graduates were employed six
months after finishing their courses. The figure attracted praise from
prospective students.

{[}G1{]}

A separate survey asked graduates whether their jobs used the skills
they had studied and whether they expected to stay in the roles.

{[}G2{]}

The college therefore decided to publish both measures rather than treat
employment alone as proof of course quality.

\begin{itemize}
\item
  \textbf{A.} Yet the number did not explain what kinds of jobs
  graduates had taken.
\item
  \textbf{B.} No graduate had studied any subject at the college.
\item
  \textbf{C.} The results suggested a more complicated picture than the
  first figure alone.
\item
  \textbf{D.} The second survey was completed before the college was
  founded.
\end{itemize}

\subsection{C　英译汉｜信息不能丢，也不能增加}\label{c-ux82f1ux8bd1ux6c49ux4fe1ux606fux4e0dux80fdux4e22ux4e5fux4e0dux80fdux589eux52a0}

\textbf{T01.} As qualifications become more widespread, their value as a
signal of unusual ability may decline.

译文：

\rule{0.8\linewidth}{0.35pt}

\textbf{T02.} Removing a degree requirement does not mean abandoning
assessment of the skills a job demands.

译文：

\rule{0.8\linewidth}{0.35pt}

\textbf{T03.} To the extent that dismissal rules increase hiring costs,
employers may become more cautious about new appointments.

译文：

\rule{0.8\linewidth}{0.35pt}

\textbf{T04.} A policy should be judged by its observed effects as well
as by the intention behind it.

译文：

\rule{0.8\linewidth}{0.35pt}

\subsection{D　汉译英｜使用准确而自然的表达}\label{d-ux6c49ux8bd1ux82f1ux4f7fux7528ux51c6ux786eux800cux81eaux7136ux7684ux8868ux8fbe}

\textbf{P01.}
新的招聘方式扩大了申请者范围，但我们仍需检验它是否真正提高了招聘质量。

试译：

\rule{0.8\linewidth}{0.35pt}

\textbf{P02.}
程序公平有其价值，但一项制度的实际效果还取决于它适用于哪些岗位。

试译：

\rule{0.8\linewidth}{0.35pt}

\subsection{E　写作小任务｜把语言知识用于表达}\label{e-ux5199ux4f5cux5c0fux4efbux52a1ux628aux8bedux8a00ux77e5ux8bc6ux7528ux4e8eux8868ux8fbe}

\textbf{任务（英文）}　Write 90--120 words for a student careers
newsletter. Discuss whether employers should always require a degree.
Give one benefit of qualifications, one limitation, and a practical way
to judge applicants.

\textbf{完成要求}　先拟两到三个论点，每段一句主题句。必须有一个明确的限定（some
/ may / under certain
conditions等），并说明例子支持的范围。写完自查：是否把一种可能夸大为必然？

\textbf{自己的草稿}

\vspace{2cm}\rule{0.95\linewidth}{0.35pt}\vspace{0.6cm}

\clearpage

\section{答案、证据与思维诊断}\label{ux7b54ux6848ux8bc1ux636eux4e0eux601dux7ef4ux8bcaux65ad}

这里既不是``答案字母库''，也不是``命题者的秘密''。重点是逐步验证自己的理解：\textbf{词义成立吗、原文说了吗、题干要求这一层信息吗、错误项改变了什么？}

\subsection{一　开始前的四道判断}\label{ux4e00-ux5f00ux59cbux524dux7684ux56dbux9053ux5224ux65ad}

\textbf{诊断 01}　机会的开启不等于对长期结果的保证。

\textbf{诊断 02}　be found wanting
表示承诺经检验未达要求，wanting不是想要。

\textbf{诊断 03}　intended to protect 属意图；costs fell
unevenly是分配效果。

\textbf{诊断 04}　to the extent that
缩小推断适用范围，不能推广所有雇主。

\subsection{二　原卷段落回访｜按证据回答，不按印象记忆}\label{ux4e8c-ux539fux5377ux6bb5ux843dux56deux8bbfux6309ux8bc1ux636eux56deux7b54ux4e0dux6309ux5370ux8c61ux8bb0ux5fc6}

\subsubsection{单元一｜大学学位：一种承诺何时不再充分}\label{ux5355ux5143ux4e00ux5927ux5b66ux5b66ux4f4dux4e00ux79cdux627fux8bfaux4f55ux65f6ux4e0dux518dux5145ux5206-1}

\textbf{D1}　首段提出重新评估选择与价值，而非普遍否定大学。26题关注reassess
necessity，不是延迟申请。

\textbf{D2}　承诺是``获得体面学位就能终身得到保障''；数字是支撑该承诺不再稳固的论据。27题要求从数据回到文章判断。

\textbf{D3}　作者没有宣称大学教育毫无价值；他限制了``必须上大学''的普遍性。

\textbf{D4}　不是，删除某些岗位的学历门槛不意味着不要求技能。28题围绕对学校毕业生的务实评价。

\textbf{D5}　对把学位作为理想路径的人提出进一步知识与技能深化，29题聚焦specialised
studies。

\textbf{D6}　作者认为新一代的学位未必如旧一代那样固定身份；30题对应career-long
learning，而不是缩窄工作选择。

\subsubsection{单元二｜劳动保障：目的与后果之间有几层推理}\label{ux5355ux5143ux4e8cux52b3ux52a8ux4fddux969cux76eeux7684ux4e0eux540eux679cux4e4bux95f4ux6709ux51e0ux5c42ux63a8ux7406-1}

\textbf{E1}　两者均是要求：show cause和act in a procedurally fair
way。36题答案是保护普通劳动者权益。

\textbf{E2}　作者讨论高薪、表现不佳的管理层，不是所有劳动者或一切保护规定。

\textbf{E3}　限制某些低绩效经理的解雇被作者视为影响管理质量和效率，进而可能危及工作机会；这是作者的分析链。

\textbf{E4}　不是；条件和may都收窄了结论。真题39题为非官方最匹配项，与``所有现有工资立即下调''不同。

\textbf{E5}　报告关联管理质量与生产率；法律保护如何影响全部经济后果则是作者延展的政策判断。

\textbf{E6}　它定义制度是否适用于某个收入群体，而不是设定员工工资必须达到的最低水平。

\subsection{三　历年真题的参考答案与选项比较}\label{ux4e09-ux5386ux5e74ux771fux9898ux7684ux53c2ux8003ux7b54ux6848ux4e0eux9009ux9879ux6bd4ux8f83}

以下是以试卷文字及独立转录交叉核对的\textbf{非官方参考答案}。涉及现实政策及历史研究时，只说明文章自己的论证；不把作者立场当作已证实的普遍规律。

\subsubsection{2022 年英语（一） Reading Text
2（26---30题）}\label{ux5e74ux82f1ux8bedux4e00-reading-text-22630ux9898}

\textbf{第26题　C　｜作者建议}

\textbf{为何选择}：首段邀请新一代重新考虑学位意义和必要性，非要求推迟申请。

\textbf{其余选项的偏差}：A仅选择学校、B努力程度、D拖延申请都没有抓住重新评估高等教育必要性的焦点。

\textbf{第27题　B　｜论据作用}

\textbf{为何选择}：P2就业比例支持``学位普及后相对区分价值减弱''的判断。

\textbf{其余选项的偏差}：A代际观点、C公众不满未由数字显示；D社会流动路径是附近话题，但没有保留``价值缩水''的否定方向。

\textbf{第28题　C　｜态度依据}

\textbf{为何选择}：作者认为某些雇主取消特定岗位学位要求是一种务实变化。

\textbf{其余选项的偏差}：A鼓励拿学位、B评价毕业生的愿望、D父母已经改观都与段落支撑的雇主态度不同。

\textbf{第29题　D　｜面向特定人群的建议}

\textbf{为何选择}：对选择学位路径者，作者建议进一步深化专门领域学习。

\textbf{其余选项的偏差}：A尽早锁定职业、B仅限在职培训、C与高薪研究生合作均不是该段的建议。

\textbf{第30题　A　｜末段概括}

\textbf{为何选择}：作者强调学习将贯穿多次职业变化，学位不再唯一界定身份。

\textbf{其余选项的偏差}：B成为教师、C不再喜欢学位、D职业选择有限均偏离持续学习的主线。

\subsubsection{2022 年英语（一） Reading Text
4（36---40题）}\label{ux5e74ux82f1ux8bedux4e00-reading-text-43640ux9898}

\textbf{第36题　D　｜制度目的}

\textbf{为何选择}：P1---2说明申诉条款旨在防止无理由解雇，保护普通雇员。

\textbf{其余选项的偏差}：A处罚公司、B改善招聘、C免除雇主责任不是法律设计目标。

\textbf{第37题　A　｜作者观点范围}

\textbf{为何选择}：P3声称限制解雇部分低绩效高薪经理可能妨碍企业改善绩效；这是文章的政策判断。

\textbf{其余选项的偏差}：B经理权威、C公众形象、D劳资关系均未触及论述的生产率机制。

\textbf{第38题　D　｜间接推断}

\textbf{为何选择}：文章引用生产率委员会将管理能力弱列为原因，据此最接近支持改善低绩效管理。

\textbf{其余选项的偏差}：A压工资、B加强一切保护、C限制企业主权限都不针对管理能力这一论据。注意委员会并未逐字提出D项完整措施，是题内合理推断。

\textbf{第39题　B　｜机制和措辞精度}

\textbf{为何选择}：作者认为解雇程序增加成本，企业可能降低薪酬水平；在所给四项中B最吻合。

\textbf{其余选项的偏差}：A高薪经理失业、C整体福祉增长、D必须增聘员工均未对应。注意lower
wages不能自动解读成所有现有员工立即名义减薪。

\textbf{第40题　C　｜跨国例子作用}

\textbf{为何选择}：澳大利亚将超过收入门槛的雇员排除在某项保护外，作者以此论述减少企业负担的可能。

\textbf{其余选项的偏差}：A保证收入、B已出现不良结果、D机制难以实施均不是对澳大利亚门槛做法的证据。

\subsection{四　句法工坊的结构与译文}\label{ux56db-ux53e5ux6cd5ux5de5ux574aux7684ux7ed3ux6784ux4e0eux8bd1ux6587}

\textbf{S1}　that限定promise的内容；as表示伴随变化，单句不强制证明唯一因果。

参考译文：随着更多毕业生进入劳动力市场，``学位能确保就业''的承诺受到质疑。

\textbf{S2}　否定极端判断，转向充分性这一更窄的问题。

参考译文：这并不是声称资格毫无用处，而是追问资格是否已经足够。

\textbf{S3}　while在此表示同时兼让步：学历门槛放松不等于能力门槛消失。

参考译文：雇主可以取消学位要求，同时仍要求提供实践能力的证据。

\textbf{S4}　前者是意图，后者是特定适用条件下的效果，不应混为一谈。

参考译文：虽然这项规则旨在保护劳动者，其效果仍会因适用岗位不同而变化。

\textbf{S5}　to the extent that + may
两重限定，使条件成立与结论概率均非绝对。

参考译文：在该程序确实增加成本的程度上，企业可能会更加谨慎地招聘。

\textbf{S6}　证据支持管理质量与生产率的关联；全部法律对工资的影响需要额外证明。

参考译文：一份将管理质量与生产率联系起来的报告，本身不能证明每项解雇规则都会降低工资。

\subsection{五　原创阅读：答案、证据与干扰分析}\label{ux4e94-ux539fux521bux9605ux8bfbux7b54ux6848ux8bc1ux636eux4e0eux5e72ux6270ux5206ux6790}

\subsubsection{Reading A · The Hiring Test That Changed the
Queue}\label{reading-a-the-hiring-test-that-changed-the-queue-1}

\textbf{A1　A}　P1把学历要求替换成实践任务；没有自动录用或强制第二学位。

\textbf{A2　B}　P2将数量反应与招聘质量分开，申请人数不等于岗位适配质量。

\textbf{A3　C}　P3是设备条件与时间差异；未提图书馆关闭。

\textbf{A4　D}　P3明确提供设备并延长期限，没有停止测试。

\textbf{A5　A}　全文区分``扩大入口''与``长期招募质量''，不作极端断言。

\textbf{AO1}　观察期短，尚需比较长期留任率和实际工作表现。

\textbf{AO2}　减少学历门槛，但可能提高设备与时间资源方面的门槛。

\textbf{本文论证复盘}：先找P1问题提出，再看中段出现的相反考量，最后将末段的限定结论写成一句``不超过原文证据范围''的话。

\subsubsection{Reading B · What a Fair Complaint Procedure Can and
Cannot
Do}\label{reading-b-what-a-fair-complaint-procedure-can-and-cannot-do-1}

\textbf{B1　B}　P1制度目的为提供解释与独立复核，并非保证申诉结果。

\textbf{B2　A}　P2提供多个竞争解释，不能以总数断定恶化或改善。

\textbf{B3　C}　P3对比即时文书成本与后续争议减少的潜在收益。

\textbf{B4　D}　P4列前三项及重复投诉，无国际市场份额。

\textbf{B5　B}　P4要求平衡质疑权和日常运作成本，反对单指标简化判断。

\textbf{BO1}　增长可能反映更多真实问题，也可能反映报告意愿增强。

\textbf{BO2}　意图看正式设计的保障目的，效果需有多指标与时间跨度的观察。

\textbf{本文论证复盘}：先找P1问题提出，再看中段出现的相反考量，最后将末段的限定结论写成一句``不超过原文证据范围''的话。

\subsection{六　专项练习答案}\label{ux516d-ux4e13ux9879ux7ec3ux4e60ux7b54ux6848}

\subsubsection{选词填空
V01---V08}\label{ux9009ux8bcdux586bux7a7a-v01v08}

\begin{itemize}
\item
  \textbf{V01}　weigh
  up。注意用法是否符合该句语法，不能只凭中文近义判断。
\item
  \textbf{V02}　found
  wanting。注意用法是否符合该句语法，不能只凭中文近义判断。
\item
  \textbf{V03}　scrap。注意用法是否符合该句语法，不能只凭中文近义判断。
\item
  \textbf{V04}　upskill。注意用法是否符合该句语法，不能只凭中文近义判断。
\item
  \textbf{V05}　grievance。注意用法是否符合该句语法，不能只凭中文近义判断。
\item
  \textbf{V06}　arbitrary。注意用法是否符合该句语法，不能只凭中文近义判断。
\item
  \textbf{V07}　threshold。注意用法是否符合该句语法，不能只凭中文近义判断。
\item
  \textbf{V08}　unwieldy。注意用法是否符合该句语法，不能只凭中文近义判断。
\end{itemize}

\subsubsection{篇章衔接 G1---G2}\label{ux7bc7ux7ae0ux8854ux63a5-g1g2}

\begin{itemize}
\item
  \textbf{G1}　A。G1承接就业率数字、指出信息缺口，为后文补充技能匹配调查铺垫。
\item
  \textbf{G2}　C。G2从两项调查比较转向更全面评价，顺势引出双指标公开的决策。
\end{itemize}

\subsubsection{英译汉 T01---T04}\label{ux82f1ux8bd1ux6c49-t01t04}

\textbf{T01}　随着资格越来越普遍，它们作为特殊能力信号的价值可能下降。

\textbf{观察要点}　as伴随变化，signal意义比``绝对价值''窄。

\textbf{T02}　取消学位要求并不意味着放弃评估岗位所需的技能。

\textbf{观察要点}　not mean + -ing；避免错误推断。

\textbf{T03}　在解雇规则确实增加招聘成本的程度上，雇主可能对新聘用更加谨慎。

\textbf{观察要点}　to the extent that与may双重限定。

\textbf{T04}　评价一项政策，不仅要看背后的初衷，也要看已观察到的效果。

\textbf{观察要点}　as well as并列，intention与effects区分。

\subsubsection{汉译英 P01---P02}\label{ux6c49ux8bd1ux82f1-p01p02}

\textbf{P01}　The new hiring approach widened access, but we still need
to examine whether it improved recruitment quality.

\textbf{观察要点}　widened access不等于improved quality。

\textbf{P02}　Procedural fairness matters, but the effects of a rule
also depend on the positions to which it applies.

\textbf{观察要点}　effect范围被positions限制。

\subsubsection{写作范文｜不要照抄，先判断它如何满足任务}\label{ux5199ux4f5cux8303ux6587ux4e0dux8981ux7167ux6284ux5148ux5224ux65adux5b83ux5982ux4f55ux6ee1ux8db3ux4efbux52a1}

\textbf{原创参考短文（约100词）}

A degree can provide valuable knowledge and show that a student has
completed a demanding programme. For some professions, formal training
is essential. However, a degree alone does not prove that a person can
perform every task in a workplace. Employers should therefore decide
which positions genuinely need specific qualifications. For other roles,
a carefully designed practical assessment may offer useful additional
evidence. Such tests should be accessible to applicants with different
resources. The goal is not to dismiss university education, but to
evaluate skills fairly while preserving the value of specialised study.
This approach offers more than a choice between degrees and no degrees.

\textbf{自己校改的三个问题}　是否明确承认学历的真实价值？是否限制了``某些岗位/另一些岗位''？实践考核是否兼顾公平可及？

\textbf{复习建议}　首次订正时用已有线索重做错题；过几天只看英文新句，不看解析重述意思。间隔时间可随真实掌握情况调整，不能把``讲过''当``会了''。

\clearpage

\section{英中双语对照附录}\label{ux82f1ux4e2dux53ccux8bedux5bf9ux7167ux9644ux5f55}

本附录集中收录本册\textbf{自编英文}材料的中译，帮助做题后逐句核对。真题整篇原文及四选项请使用前述试卷链接查阅，本册不大量转录。

\subsection{A　入门诊断｜英中逐句}\label{a-ux5165ux95e8ux8bcaux65adux82f1ux4e2dux9010ux53e5}

\textbf{诊断 01}　A qualification can open a door without guaranteeing a
career.

\textbf{中译}　一项资格可以打开机会之门，却不能保证整个职业生涯。

\textbf{诊断 02}　The promise was found wanting when graduates entered
the job market.

\textbf{中译}　当毕业生进入劳动力市场时，这一承诺被证明不尽如人意。

\textbf{诊断 03}　The policy was intended to protect employees, yet its
costs fell unevenly.

\textbf{中译}　该政策旨在保护雇员，但成本分布并不均衡。

\textbf{诊断 04}　To the extent that the rule delays hiring, new
applicants may be affected.

\textbf{中译}　在该规则确实延缓招聘的范围内，新申请者可能受到影响。

\subsection{B　单元一｜大学学位：一种承诺何时不再充分｜原创新句}\label{b-ux5355ux5143ux4e00ux5927ux5b66ux5b66ux4f4dux4e00ux79cdux627fux8bfaux4f55ux65f6ux4e0dux518dux5145ux5206ux539fux521bux65b0ux53e5}

\textbf{D1}　Before choosing a course, applicants should weigh up both
costs and possible benefits.

\textbf{中译}　选择课程前，申请者应权衡成本和可能收益。

\textbf{D2}　The training was useful, but the promise of automatic
promotion was found wanting.

\textbf{中译}　培训有用，但自动升职的承诺被证明未达到预期。

\textbf{D3}　This is not to deny the value of university, but to
question whether one route suits everyone.

\textbf{中译}　这并非否定大学的价值，而是质疑同一条路径是否适合所有人。

\textbf{D4}　Some employers removed degree requirements for roles that
relied on practical skills.

\textbf{中译}　一些雇主取消了主要依靠实践技能岗位的学位要求。

\textbf{D5}　A general education can be strengthened by specialised
training in a chosen field.

\textbf{中译}　通识教育可以通过所选领域的专门训练得到加强。

\textbf{D6}　A degree may begin a career trajectory without defining the
whole journey.

\textbf{中译}　学位可以开启职业轨迹，却不必定义整个职业历程。

\subsubsection{词汇例句的翻译}\label{ux8bcdux6c47ux4f8bux53e5ux7684ux7ffbux8bd1}

\textbf{weigh up}　Applicants must weigh up tuition costs against
opportunities for practical training.

\textbf{中译}　申请者必须权衡学费与获得实践培训的机会。

\textbf{be found wanting}　The claim that one certificate guaranteed a
job was eventually found wanting.

\textbf{中译}　一张证书就能保证就业的说法最终被证明不充分。

\textbf{devalue}　When nearly everyone held the badge, it no longer
distinguished applicants and became devalued.

\textbf{中译}　当几乎人人都有这一徽章时，它不再能区分申请者，因此相对贬值。

\textbf{social mobility}　Affordable training can offer a route to
social mobility for some families.

\textbf{中译}　负担得起的培训可能为一些家庭提供社会流动的路径。

\textbf{inevitable}　The move from secondary school to university is not
inevitable for every learner.

\textbf{中译}　从中学升入大学并非每位学习者都必然走过的路径。

\textbf{predecessor}　Young engineers often learn from the mistakes of
their predecessors.

\textbf{中译}　年轻工程师常从前辈的失误中学习。

\textbf{committed}　A committed trainee can make progress without
following a conventional degree route.

\textbf{中译}　一位投入的学员即使不走传统学位路线，也能取得进步。

\textbf{scrap}　The employer scrapped a degree requirement but retained
the skills test.

\textbf{中译}　雇主取消了学位要求，但保留了技能测试。

\textbf{generalist}　A generalist may connect ideas across departments
but still need deep expertise in one task.

\textbf{中译}　通才可以联系不同部门的想法，但仍可能需要对一项任务有深入专长。

\textbf{postgraduate}　Postgraduate study may deepen expertise, although
its value varies by field.

\textbf{中译}　研究生学习可能加深专长，但其价值因领域而异。

\textbf{upskill}　Workers can upskill through short courses while
remaining in employment.

\textbf{中译}　劳动者可以在职期间通过短课程提升技能。

\textbf{employable}　Technical knowledge and communication skills can
help graduates remain employable.

\textbf{中译}　专业知识与沟通能力能帮助毕业生保持就业竞争力。

\textbf{trajectory}　Her career trajectory included teaching, design and
later project management.

\textbf{中译}　她的职业轨迹包括教学、设计及后来的项目管理。

\textbf{fulfillment}　Some people value personal fulfillment more than a
prestigious job title.

\textbf{中译}　有些人看重个人满足感，胜过体面的工作头衔。

\subsection{B　单元二｜劳动保障：目的与后果之间有几层推理｜原创新句}\label{b-ux5355ux5143ux4e8cux52b3ux52a8ux4fddux969cux76eeux7684ux4e0eux540eux679cux4e4bux95f4ux6709ux51e0ux5c42ux63a8ux7406ux539fux521bux65b0ux53e5}

\textbf{E1}　A dismissal rule may require an employer to explain the
reason and follow a fair procedure.

\textbf{中译}　一项解雇规则可能要求雇主说明理由并遵循公平程序。

\textbf{E2}　A rule designed for new workers may have different effects
when applied to senior executives.

\textbf{中译}　原为新员工设计的规则，用于高层管理者时可能产生不同效果。

\textbf{E3}　A policy can be well intended yet produce costs through the
choices it changes.

\textbf{中译}　一项政策可能初衷良好，但会通过它改变的行为选择产生成本。

\textbf{E4}　To the extent that hiring becomes more costly, a firm may
recruit more cautiously.

\textbf{中译}　在招聘确实变得更昂贵的程度上，企业可能更加谨慎地招人。

\textbf{E5}　The report identified weak management as one possible
source of low productivity.

\textbf{中译}　报告将管理薄弱列为生产率低下的一个可能来源。

\textbf{E6}　An income threshold excludes some senior staff from a rule
while leaving others covered.

\textbf{中译}　收入门槛把部分高层雇员排除在某项规则之外，同时使其他人继续受到覆盖。

\subsubsection{词汇例句的翻译}\label{ux8bcdux6c47ux4f8bux53e5ux7684ux7ffbux8bd1-1}

\textbf{grievance}　The employee filed a grievance after receiving no
explanation for the dismissal.

\textbf{中译}　员工因遭解雇时未获解释而提出申诉。

\textbf{provision}　The provision requires employers to keep a written
record of the decision.

\textbf{中译}　该条款要求雇主保留决定的书面记录。

\textbf{dismissal}　A dismissal can be lawful only if the required
procedures are followed.

\textbf{中译}　只有遵守必要程序，解雇才可能合法。

\textbf{justify}　The manager could not justify the decision with
reliable evidence.

\textbf{中译}　经理无法用可靠证据证明该决定正当。

\textbf{procedural}　Procedural fairness requires both notice and a
chance to respond.

\textbf{中译}　程序公平要求事先告知并给予回应机会。

\textbf{safeguard}　A transparent review process can safeguard employees
against arbitrary decisions.

\textbf{中译}　透明的复核程序能够保障员工免受任意决定的伤害。

\textbf{arbitrary}　The new procedure aims to prevent arbitrary changes
to pay.

\textbf{中译}　新程序旨在防止随意更改薪酬。

\textbf{constrain}　High training costs may constrain the number of
apprentices a firm can hire.

\textbf{中译}　高昂培训成本可能限制企业能招聘的学徒数量。

\textbf{productivity}　The factory improved productivity without
increasing the number of working hours.

\textbf{中译}　工厂没有增加工时，却提高了生产率。

\textbf{paradoxically}　Paradoxically, a rule meant to speed up hiring
created more paperwork.

\textbf{中译}　颇为反常的是，一项旨在加快招聘的规则却增加了文书工作。

\textbf{marginal}　A marginal applicant may be excluded when recruitment
costs rise.

\textbf{中译}　当招聘成本上升时，处于边缘位置的申请者可能被排除。

\textbf{stringent}　Stringent reporting rules may improve transparency
but raise administrative costs.

\textbf{中译}　严格的报告规则可能提高透明度，但会增加管理成本。

\textbf{threshold}　Employees above the income threshold are covered by
a different review procedure.

\textbf{中译}　收入超过门槛的员工适用另一套复核程序。

\textbf{unwieldy}　The proposed appeal system was too unwieldy for small
employers to use.

\textbf{中译}　拟议的申诉制度对小雇主而言过于繁琐，难以使用。

\subsection{C　句法工坊｜六句双语}\label{c-ux53e5ux6cd5ux5de5ux574aux516dux53e5ux53ccux8bed}

\textbf{S1}　The promise that a degree would secure a job was questioned
as more graduates entered the labour market.

\textbf{中译}　随着更多毕业生进入劳动力市场，``学位能确保就业''的承诺受到质疑。

\textbf{S2}　This is not to claim that qualifications are useless, but
to ask whether they are sufficient.

\textbf{中译}　这并不是声称资格毫无用处，而是追问资格是否已经足够。

\textbf{S3}　An employer may remove a degree requirement while still
demanding evidence of practical competence.

\textbf{中译}　雇主可以取消学位要求，同时仍要求提供实践能力的证据。

\textbf{S4}　Although the rule was intended to protect workers, its
effects varied with the positions to which it applied.

\textbf{中译}　虽然这项规则旨在保护劳动者，其效果仍会因适用岗位不同而变化。

\textbf{S5}　To the extent that the procedure adds costs, firms may
become more cautious about hiring.

\textbf{中译}　在该程序确实增加成本的程度上，企业可能会更加谨慎地招聘。

\textbf{S6}　A report linking management quality to productivity does
not by itself prove that every dismissal rule reduces wages.

\textbf{中译}　一份将管理质量与生产率联系起来的报告，本身不能证明每项解雇规则都会降低工资。

\subsection{D1　Reading A · The Hiring Test That Changed the
Queue｜全文双语}\label{d1-reading-a-the-hiring-test-that-changed-the-queueux5168ux6587ux53ccux8bed}

\textbf{P1（英文）}　When a design company removed the degree
requirement from its entry-level vacancies, the change attracted
attention. Managers explained that many successful employees had learned
their skills outside university. The new advertisements asked applicants
to complete a practical task instead of listing their academic
qualifications. The firm described the policy as a way to widen access,
but it offered no guarantee that everyone would have an equal chance of
being hired.

\textbf{P1（中译）}　一家设计公司取消初级岗位的学位要求后，这一变化引起关注。经理解释称，许多优秀员工是在大学之外掌握技能的。新招聘广告要求申请者完成实践任务，而不是列出学历。公司把这项政策描述为扩大机会的方式，但并未保证人人都有平等录用机会。

\textbf{P2（英文）}　Applications rose sharply during the first month.
Some observers took this as proof that the scheme had succeeded. Yet an
increase in applications was an early response, not evidence about the
quality of recruitment. The company still needed to know whether the
tests identified people who could do the job and whether the extra
assessment time was reasonable.

\textbf{P2（中译）}　第一个月申请数量显著增加。一些观察者将此视为计划成功的证明。然而，申请数量增加只是早期反应，并不能证明招聘质量。公司仍需知道测试能否识别能胜任工作的人，以及额外测评时间是否合理。

\textbf{P3（英文）}　The practical test also created a new obstacle.
Applicants with access to expensive equipment could complete it quickly,
while others had to borrow a computer or find a quiet room. The absence
of a degree requirement had removed one barrier but left others in
place. After receiving this feedback, the firm began offering equipment
at local libraries and extended the submission deadline.

\textbf{P3（中译）}　实践测试也造成新障碍。能使用昂贵设备的申请者可以很快完成，而其他人必须借电脑或寻找安静空间。取消学位要求移除了一道壁垒，但仍留有其他壁垒。收到反馈后，公司开始在地方图书馆提供设备，并延长提交期限。

\textbf{P4（英文）}　Six months later, the firm reported that applicants
from a wider range of backgrounds had reached the interview stage. It
did not claim that the new approach had increased long-term
productivity. Managers planned to compare retention rates and
performance over several years before drawing a stronger conclusion. The
example suggests that opening an entrance is a useful first step, but it
does not tell us everything about what happens after people enter.

\textbf{P4（中译）}　六个月后，公司报告称，进入面试阶段的申请者背景更为多样。但公司没有声称新做法已经提高了长期生产率。管理人员计划在数年间比较员工留任率和表现，再得出更强的结论。这个例子说明，打开入口是有用的第一步，却不能告诉我们人们进入后发生的一切。

\subsubsection{阅读题及全部选项英中对照}\label{ux9605ux8bfbux9898ux53caux5168ux90e8ux9009ux9879ux82f1ux4e2dux5bf9ux7167}

\textbf{A1.}　What replaced the degree requirement?

\textbf{中译}　什么取代了学位要求？

\begin{itemize}
\item
  \textbf{A.} A practical task.

  \textbf{中译}　一项实践任务。
\item
  \textbf{B.} A compulsory second degree.

  \textbf{中译}　强制要求第二个学位。
\item
  \textbf{C.} An unpaid two-year contract.

  \textbf{中译}　一份无薪两年合同。
\item
  \textbf{D.} An automatic interview.

  \textbf{中译}　自动获得面试。
\end{itemize}

\textbf{A2.}　Why was the rise in applications insufficient evidence of
success?

\textbf{中译}　申请增加为什么不足以证明成功？

\begin{itemize}
\item
  \textbf{A.} It occurred several years later.

  \textbf{中译}　它发生在数年之后。
\item
  \textbf{B.} It said little about recruitment quality.

  \textbf{中译}　它无法充分说明招聘质量。
\item
  \textbf{C.} It was limited to university graduates.

  \textbf{中译}　它仅涉及大学毕业生。
\item
  \textbf{D.} It proved the test unfair.

  \textbf{中译}　它已证明测试不公。
\end{itemize}

\textbf{A3.}　Some applicants were disadvantaged because

\textbf{中译}　部分申请者处于不利地位，是因为？

\begin{itemize}
\item
  \textbf{A.} the company refused all applications.

  \textbf{中译}　公司拒绝全部申请。
\item
  \textbf{B.} the test was about economic theory.

  \textbf{中译}　测试主题是经济理论。
\item
  \textbf{C.} their access to equipment was limited.

  \textbf{中译}　他们获取设备的条件受限。
\item
  \textbf{D.} libraries had been permanently closed.

  \textbf{中译}　图书馆被永久关闭。
\end{itemize}

\textbf{A4.}　The firm changed its procedure by

\textbf{中译}　公司怎样调整程序？

\begin{itemize}
\item
  \textbf{A.} removing the practical test entirely.

  \textbf{中译}　完全取消实践测试。
\item
  \textbf{B.} reintroducing the degree rule.

  \textbf{中译}　重新引入学位要求。
\item
  \textbf{C.} promising every applicant a job.

  \textbf{中译}　承诺每个申请者都有工作。
\item
  \textbf{D.} offering equipment and more time.

  \textbf{中译}　提供设备和更多时间。
\end{itemize}

\textbf{A5.}　The author's central point is that

\textbf{中译}　作者核心观点是？

\begin{itemize}
\item
  \textbf{A.} wider access is only one part of evaluating reform.

  \textbf{中译}　扩大进入机会只是评价改革的一部分。
\item
  \textbf{B.} all academic qualifications are worthless.

  \textbf{中译}　所有学历资格都没有价值。
\item
  \textbf{C.} practical tests always reduce fairness.

  \textbf{中译}　实践测试总会降低公平。
\item
  \textbf{D.} early application numbers prove productivity gains.

  \textbf{中译}　初期申请数证明生产率提高。
\end{itemize}

\subsubsection{无选项简答的题干对照}\label{ux65e0ux9009ux9879ux7b80ux7b54ux7684ux9898ux5e72ux5bf9ux7167}

\textbf{O1}　Why did the firm avoid claiming higher productivity?

\textbf{中译}　公司为什么不声称生产率已经提高？

\textbf{O2}　What trade-off did the practical test introduce?

\textbf{中译}　实践测试带来了什么新的取舍？

\subsection{D2　Reading B · What a Fair Complaint Procedure Can and
Cannot
Do｜全文双语}\label{d2-reading-b-what-a-fair-complaint-procedure-can-and-cannot-doux5168ux6587ux53ccux8bed}

\textbf{P1（英文）}　A medium-sized logistics firm introduced a new
complaint procedure after several employees said that decisions about
shifts were difficult to challenge. Under the procedure, managers had to
explain changes in writing, and employees could request a review by
someone who had not made the original decision. The policy was meant to
reduce arbitrary treatment, not to prevent the company from changing
schedules whenever necessary.

\textbf{P1（中译）}　一家中型物流企业在数位员工表示排班决定难以申诉之后，引入新的投诉程序。根据该程序，经理必须书面解释变更，员工可以要求由未作出原决定的人复核。政策旨在减少任意对待，而不是阻止公司在有必要时调整排班。

\textbf{P2（英文）}　At first, the number of formal complaints
increased. Some managers viewed this as a sign that the workplace had
become less harmonious. An employee representative offered another
interpretation: workers may simply have become more willing to report
problems that previously went unrecorded. Neither interpretation could
be established from the complaint count alone.

\textbf{P2（中译）}　最初，正式投诉数量增加。一些经理把这看作工作环境变得不和谐的信号。一位员工代表提出另一种解释：劳动者可能只是更愿意报告过去没有被记录的问题。仅凭投诉数量，无法证实哪种解释成立。

\textbf{P3（英文）}　The procedure also required time. Supervisors said
that documenting small changes left them less time for daily operations.
The employee representative argued that clear records could prevent
larger disputes later. The disagreement was partly about when the costs
and benefits would appear. A report that examined only the first month
would capture paperwork, but not necessarily avoided disputes six months
later.

\textbf{P3（中译）}　该程序也耗费时间。主管表示，记录小调整让他们用于日常运作的时间减少。员工代表认为，清楚记录可防止之后更大的争议。分歧部分在于成本和收益何时出现。只研究第一个月的报告能看到文书成本，却未必能观察到六个月后避免的争议。

\textbf{P4（英文）}　The firm therefore agreed to review the policy
after a full year, using several measures: time spent on paperwork,
repeat complaints, staff turnover and employees' confidence in the
process. No single number would settle whether the procedure was fair or
efficient. The more useful question was whether it could protect the
right to challenge a decision without making routine work unnecessarily
difficult.

\textbf{P4（中译）}　因此公司同意在完整一年后，以多项指标审查政策，包括文书时间、重复投诉、员工流动率及员工对流程的信心。没有单一数字能断定该程序公平或高效。更有价值的问题是，它能否保护员工质疑决定的权利，同时不让日常工作变得不必要地困难。

\subsubsection{阅读题及全部选项英中对照}\label{ux9605ux8bfbux9898ux53caux5168ux90e8ux9009ux9879ux82f1ux4e2dux5bf9ux7167-1}

\textbf{B1.}　The new procedure was mainly designed to

\textbf{中译}　新程序主要旨在？

\begin{itemize}
\item
  \textbf{A.} abolish all schedule changes.

  \textbf{中译}　废除所有排班变动。
\item
  \textbf{B.} give workers a way to challenge decisions.

  \textbf{中译}　给劳动者提供质疑决定的渠道。
\item
  \textbf{C.} guarantee every complaint would succeed.

  \textbf{中译}　保证每项投诉成功。
\item
  \textbf{D.} replace all supervisors.

  \textbf{中译}　取代所有主管。
\end{itemize}

\textbf{B2.}　The rise in complaints could mean that

\textbf{中译}　投诉增加可能意味着？

\begin{itemize}
\item
  \textbf{A.} workers had become more willing to report problems.

  \textbf{中译}　劳动者更愿意报告问题。
\item
  \textbf{B.} all managers were acting illegally.

  \textbf{中译}　所有经理都违法。
\item
  \textbf{C.} the procedure had certainly failed.

  \textbf{中译}　程序肯定失败。
\item
  \textbf{D.} no workplace problems existed before.

  \textbf{中译}　以前没有职场问题。
\end{itemize}

\textbf{B3.}　What did the two sides disagree about in Paragraph 3?

\textbf{中译}　第三段双方主要争论什么？

\begin{itemize}
\item
  \textbf{A.} Whether any record should exist at all.

  \textbf{中译}　是否应完全不作记录。
\item
  \textbf{B.} Who should own the company.

  \textbf{中译}　谁应拥有公司。
\item
  \textbf{C.} When the procedure's costs and benefits might emerge.

  \textbf{中译}　程序成本与收益何时体现。
\item
  \textbf{D.} Whether shifts were legal in principle.

  \textbf{中译}　排班原则上是否合法。
\end{itemize}

\textbf{B4.}　The annual review will measure all except

\textbf{中译}　年度审查将衡量除哪项以外的指标？

\begin{itemize}
\item
  \textbf{A.} staff turnover.

  \textbf{中译}　员工流动率。
\item
  \textbf{B.} time spent on documentation.

  \textbf{中译}　文书时间。
\item
  \textbf{C.} trust in the procedure.

  \textbf{中译}　对程序的信任。
\item
  \textbf{D.} the firm's international market share.

  \textbf{中译}　公司的国际市场份额。
\end{itemize}

\textbf{B5.}　What is the author's preferred approach?

\textbf{中译}　作者更倾向怎样评价该制度？

\begin{itemize}
\item
  \textbf{A.} Reject it because complaints increased.

  \textbf{中译}　因投诉增加就将它否定。
\item
  \textbf{B.} Assess both procedural protection and practical costs.

  \textbf{中译}　同时考察程序保障和实际成本。
\item
  \textbf{C.} Judge success only by the number of complaints.

  \textbf{中译}　只按投诉数量判成功。
\item
  \textbf{D.} Assume all employee reports are inaccurate.

  \textbf{中译}　假定所有雇员报告都不准确。
\end{itemize}

\subsubsection{无选项简答的题干对照}\label{ux65e0ux9009ux9879ux7b80ux7b54ux7684ux9898ux5e72ux5bf9ux7167-1}

\textbf{O1}　Why can complaint numbers alone mislead?

\textbf{中译}　为什么单凭投诉数量会误导？

\textbf{O2}　How would you distinguish a policy's intention from its
result?

\textbf{中译}　怎样区分政策意图和政策效果？

\subsection{E　专项练习英中对照}\label{e-ux4e13ux9879ux7ec3ux4e60ux82f1ux4e2dux5bf9ux7167}

\subsubsection{完成后的词汇句子}\label{ux5b8cux6210ux540eux7684ux8bcdux6c47ux53e5ux5b50}

\textbf{V01}　Young applicants should weigh up the value of formal study
against other options.

\textbf{中译}　年轻申请者应将正规学习的价值与其他选择进行权衡。

\textbf{V02}　The claim of guaranteed employment was found wanting after
several years of weak results.

\textbf{中译}　在数年效果不佳之后，``保证就业''的主张被证明不足。

\textbf{V03}　The firm decided to scrap its degree requirement for one
technical role.

\textbf{中译}　公司决定取消某一技术岗位的学位要求。

\textbf{V04}　Employees can upskill their skills through short
programmes.

\textbf{中译}　雇员可以通过短期课程提升技能。

\textbf{V05}　The new procedure gives employees a way to file a
grievance.

\textbf{中译}　新程序让雇员有机会提出申诉。

\textbf{V06}　A flexible rule can help prevent arbitrary decisions by
management.

\textbf{中译}　灵活规则可以帮助防止管理层作出任意决定。

\textbf{V07}　The law applies only to employees below a specified income
threshold.

\textbf{中译}　该法仅适用于收入低于指定门槛的雇员。

\textbf{V08}　The proposed approval process proved too unwieldy for
small businesses.

\textbf{中译}　拟议审批程序对小企业而言过于繁琐。

\subsubsection{篇章衔接的原句、备选句及中译}\label{ux7bc7ux7ae0ux8854ux63a5ux7684ux539fux53e5ux5907ux9009ux53e5ux53caux4e2dux8bd1}

\textbf{英文}　A college announced that 82 percent of its graduates were
employed six months after finishing their courses. The figure attracted
praise from prospective students.

\textbf{中译}　一所学院宣布，82\%的毕业生在完成课程六个月后已就业。这个数字受到未来学生的称赞。

\textbf{英文}　A separate survey asked graduates whether their jobs used
the skills they had studied and whether they expected to stay in the
roles.

\textbf{中译}　另一项调查询问毕业生：目前岗位是否用上了学习技能，以及是否打算继续留任。

\textbf{英文}　The college therefore decided to publish both measures
rather than treat employment alone as proof of course quality.

\textbf{中译}　因此学院决定同时公布两项指标，而不再把就业本身视为课程质量的充分证明。

\textbf{A}　Yet the number did not explain what kinds of jobs graduates
had taken.

\textbf{中译}　然而，该数字无法说明毕业生从事了哪种工作。

\textbf{B}　No graduate had studied any subject at the college.

\textbf{中译}　没有一个毕业生曾在学院学习任何科目。

\textbf{C}　The results suggested a more complicated picture than the
first figure alone.

\textbf{中译}　结果表明，情况比单独的首个数字所显示的更复杂。

\textbf{D}　The second survey was completed before the college was
founded.

\textbf{中译}　第二项调查在学院建立前就完成了。

\subsubsection{翻译、表达与写作任务}\label{ux7ffbux8bd1ux8868ux8fbeux4e0eux5199ux4f5cux4efbux52a1}

\textbf{T01}　As qualifications become more widespread, their value as a
signal of unusual ability may decline.

\textbf{中译}　随着资格越来越普遍，它们作为特殊能力信号的价值可能下降。

\textbf{T02}　Removing a degree requirement does not mean abandoning
assessment of the skills a job demands.

\textbf{中译}　取消学位要求并不意味着放弃评估岗位所需的技能。

\textbf{T03}　To the extent that dismissal rules increase hiring costs,
employers may become more cautious about new appointments.

\textbf{中译}　在解雇规则确实增加招聘成本的程度上，雇主可能对新聘用更加谨慎。

\textbf{T04}　A policy should be judged by its observed effects as well
as by the intention behind it.

\textbf{中译}　评价一项政策，不仅要看背后的初衷，也要看已观察到的效果。

\textbf{P01}　The new hiring approach widened access, but we still need
to examine whether it improved recruitment quality.

\textbf{中译}　新的招聘方式扩大了申请者范围，但我们仍需检验它是否真正提高了招聘质量。

\textbf{P02}　Procedural fairness matters, but the effects of a rule
also depend on the positions to which it applies.

\textbf{中译}　程序公平有其价值，但一项制度的实际效果还取决于它适用于哪些岗位。

\textbf{写作任务（英文）}　Write 90--120 words for a student careers
newsletter. Discuss whether employers should always require a degree.
Give one benefit of qualifications, one limitation, and a practical way
to judge applicants.

\textbf{中译}　为学生职业通讯写约90---120词短文，讨论雇主是否应始终要求学位，至少提到学历的一项价值、一项局限和一种实际的选人方式。

\textbf{原创写作参考（英文）}　A degree can provide valuable knowledge
and show that a student has completed a demanding programme. For some
professions, formal training is essential. However, a degree alone does
not prove that a person can perform every task in a workplace. Employers
should therefore decide which positions genuinely need specific
qualifications. For other roles, a carefully designed practical
assessment may offer useful additional evidence. Such tests should be
accessible to applicants with different resources. The goal is not to
dismiss university education, but to evaluate skills fairly while
preserving the value of specialised study. This approach offers more
than a choice between degrees and no degrees.

\textbf{中译}　学位可以提供有价值的知识，也能表明学生完成了要求严格的课程。对于某些职业，正规训练不可或缺。然而，学位本身不能证明一个人能完成职场中的每项任务。因此，雇主应判断哪些岗位真正需要特定资格；对于其他岗位，精心设计的实践测评可以提供额外证据。这类测试还应让拥有不同资源的申请者都能够参加。目标不是贬低大学教育，而是在保留专门学习价值的同时公平评价技能。这种方法比``要求学位还是完全不要学位''的二选一更合理。

\clearpage

\section{学习材料与出处}\label{ux5b66ux4e60ux6750ux6599ux4e0eux51faux5904}

本册真实试卷伴读依据相应年份英语（一）真题的公开整理版，并结合独立转录核对段落与第26---30或36---40题。参考答案仅为核对后的非官方参考答案。不同网站出现拼写、删节或翻译差异时，以试卷版对照，不把来源网页中的解析直接当成标准答案。

\begin{itemize}
\item
  2022 年英语（一） Reading Text
  2（26---30题）：\href{https://ac-wiki.org/general-skills/study/assets/2022-201.pdf}{打开试卷页面}
\item
  2022 年英语（一） Reading Text
  4（36---40题）：\href{https://ac-wiki.org/general-skills/study/assets/2022-201.pdf}{打开试卷页面}
\end{itemize}

课外词汇学习参考：I. S. P. Nation, \emph{Learning Vocabulary in Another
Language}, Cambridge University
Press（2022），词汇知识不仅包括释义，还包括形式、搭配、语法功能与使用限制。

\textbf{使用提醒}：自编短文只用于训练，不冒充命题者写作或某一年试卷原文。它们的题目故意包含不同类型的错误项，以训练根据具体证据而不是根据口诀判断。

\end{document}
